\part{The guided tour}

% Part divider text: see the note in main.tex about \@endpart.
\noindent\textbf{\large What this part covers}

\medskip
\noindent
One chapter per screen. Each opens with a photograph of the running application,
then shows the code that produced it, then explains the decisions behind that
code. The file path is printed above every listing, so you can open the same
file on your own machine and read along.

\medskip
\noindent
These chapters are meant to be read at the keyboard with the stack running. If
it is not running yet, jump forward to Part~III, start it, and come back.

\chapter{Sign in}
\label{ch:login}

\shot{01-login-empty.png}{The sign-in screen. Everything an unauthenticated
visitor can see of IncidentFlow.}

\section{What is deliberately missing}

Look at the screenshot again and notice what is \emph{not} there: there are no
demo account credentials printed on the page.

They used to be. During development it is convenient to show ``sign in as
\fpath{admin@example.com} / \fpath{password}'' right on the form. It was removed,
and the reason is written into the file so nobody adds it back:

\where{web/src/pages/LoginPage.tsx}
\begin{lstlisting}[language=tsx]
/**
 * Sign-in.
 *
 * Note what is deliberately absent: demo account credentials. They used to be
 * printed here for convenience. A password rendered on a login screen trains
 * every visitor to treat this as a toy, and there is no way for a viewer to
 * tell "seeded demo password" from "real password left in the markup". They now
 * live in the README and the setup manual, where a reader has already chosen to
 * look at project documentation.
 */
\end{lstlisting}

\note{This is a small thing that says something larger. A visitor cannot tell
the difference between a password that is safe to show and one that is not ---
so showing any password teaches them that this application is careless with
them. The credentials did not disappear; they moved somewhere a reader arrives
deliberately.}

\section{The two halves of the screen}

The layout is a split panel: brand on the left, form on the right. The left half
vanishes on a narrow screen, and that is a decision rather than a default:

\where{web/src/pages/LoginPage.tsx}
\begin{lstlisting}[language=tsx]
{/* ---------------------------------------------------------------- */}
{/* Brand panel. Hidden below lg: on a phone it would push the actual */}
{/* form below the fold, which is the one thing a login page must not */}
{/* do.                                                              */}
{/* ---------------------------------------------------------------- */}
<aside className="relative hidden overflow-hidden bg-brand-900 p-12 lg:flex lg:flex-col lg:justify-between">
\end{lstlisting}

The \fpath{hidden} and \fpath{lg:flex} classes are Tailwind: hidden by default,
becoming a flex column only at the \fpath{lg} breakpoint and above. On a phone
the form is the first thing on screen.

\section{The password field, and why it has an eye}

\shot{03-login-password-revealed.png}{The reveal toggle, shown pressed. The
field's type flips between \fpath{password} and \fpath{text}.}

\where{web/src/pages/LoginPage.tsx}
\begin{lstlisting}[language=tsx]
<input
  id="password"
  type={showPassword ? 'text' : 'password'}
  autoComplete="current-password"
  required
  value={password}
  onChange={(event) => setPassword(event.target.value)}
  className="input pr-11"
/>
<button
  type="button"
  onClick={() => setShowPassword((visible) => !visible)}
  // Typos are invisible in a masked field and the error message
  // is deliberately unhelpful, so a reveal toggle is the only
  // way to tell a wrong password from a mistyped one.
  aria-label={showPassword ? 'Hide password' : 'Show password'}
  className="absolute inset-y-0 right-0 grid w-11 place-items-center ..."
>
\end{lstlisting}

Three details in that block are worth naming.

\begin{description}[leftmargin=1.2cm,style=nextline]
\item[\fpath{type=\{showPassword ? 'text' : 'password'\}}] React re-renders the
  input with a different \fpath{type} attribute. There is no separate ``visible''
  input; it is one field that changes character.
\item[\fpath{autoComplete="current-password"}] Tells the browser and password
  managers what this field is. Without it, saved passwords do not offer
  themselves, and users blame your application.
\item[\fpath{aria-label}] What a screen reader announces for a button whose
  visible content is an icon with no text.
\end{description}

\bugbox{That \fpath{aria-label} broke every end-to-end test the day it was
added. The tests located the password field with
\fpath{getByLabel('Password')}, and Playwright's \fpath{getByLabel} matches
\emph{substrings} --- so ``Password'' also matched ``Show password'' on the new
button. Two elements matched one locator, and Playwright refuses to guess. Full
story in Chapter~\ref{ch:bugs-frontend}.}

\section{What happens when you press Sign in}

\where{web/src/pages/LoginPage.tsx}
\begin{lstlisting}[language=tsx]
async function handleSubmit(event: FormEvent) {
  event.preventDefault();
  setError(null);
  setSubmitting(true);

  try {
    await login(email, password);
  } catch (cause) {
    // The server deliberately returns the same message for an unknown email
    // and a wrong password; passing it straight through keeps the client from
    // reintroducing the enumeration leak the API was careful to avoid.
    setError(
      cause instanceof ApiError ? cause.message : 'Sign-in failed. Check your connection and try again.',
    );
  } finally {
    setSubmitting(false);
  }
}
\end{lstlisting}

\fpath{event.preventDefault()} stops the browser's default behaviour of
reloading the page when a form is submitted. Without it, React never sees the
submission --- the page simply navigates.

\section{The error message is deliberately unhelpful}

\shot{02-login-error.png}{A failed sign-in. Read the message carefully.}

The message is \emph{``Those credentials do not match our records.''} It does
not say whether the email exists.

That is not laziness. Here is the server code:

\where{api/app/Http/Controllers/Api/V1/AuthController.php}
\begin{lstlisting}[language=PHP]
$user = User::query()->where('email', $credentials['email'])->first();

/**
 * Hash::check runs even when the user does not exist would be ideal for
 * constant time; Laravel's own attempt() has the same shape. The real
 * defence against enumeration here is the identical response for both
 * cases, plus the per-email+IP rate limit on this route.
 */
if ($user === null || ! Hash::check($credentials['password'], $user->password)) {
    return $this->invalidCredentials($request);
}
\end{lstlisting}

\note{This defends against \textbf{user enumeration}. If a wrong password said
``incorrect password'' and an unknown address said ``no such user'', anyone could
discover who has an account by trying addresses --- useful for phishing, and for
knowing which passwords are worth attacking. One message for both cases removes
that signal. The cost is a slightly worse experience for someone who genuinely
mistyped their address, which is why the reveal toggle exists.}

\subsection{A disabled account is treated differently}

\where{api/app/Http/Controllers/Api/V1/AuthController.php}
\begin{lstlisting}[language=PHP]
if (! $user->is_active) {
    return response()->json([
        'error' => [
            'code' => 'account_disabled',
            'message' => 'This account has been disabled. Contact an administrator.',
            'request_id' => $request->headers->get('X-Request-Id'),
        ],
    ], Response::HTTP_FORBIDDEN);
}
\end{lstlisting}

This message \emph{is} specific, and that is consistent rather than
contradictory: to reach it you must already have supplied the correct password,
so you are not an attacker probing for valid addresses --- you are the account
holder, and telling you to contact an administrator is the useful answer.

\subsection{Every error has the same shape}

Notice the JSON: \fpath{code}, \fpath{message}, \fpath{request\_id}. Every error
the API produces has that shape, from any endpoint, formatted in one place.

\where{api/app/Exceptions/ApiExceptionRenderer.php}
\begin{lstlisting}[language=PHP]
/**
 * One error shape for the entire API.
 *
 *   { "error": { "code", "message", "details", "request_id" } }
 *
 * Central error handling is not about tidiness. It is about two properties
 * that are impossible to maintain if each controller formats its own failures:
 *
 *  - the frontend has exactly one parser, and switches on a stable `code`
 *    rather than on prose that changes when someone rewords a message;
 *  - server faults never leak internals.
 */
\end{lstlisting}

The \fpath{request\_id} matters operationally: it appears in the response
\emph{and} in the server log, so a user reporting ``it failed'' can quote a
string that leads straight to the exact log line.

\section{What sign-in actually returns}

Two things, and they are stored differently on purpose.

\begin{center}
\begin{tabular}{@{}>{\raggedright\arraybackslash}p{3.3cm}>{\raggedright\arraybackslash}p{3.0cm}p{7.6cm}@{}}
\toprule
\textbf{Token} & \textbf{Lives} & \textbf{Why there} \\
\midrule
Access token & In memory (JavaScript) & Short-lived, 15 minutes. Sent on every API call. Never written to disk, so a stolen laptop does not hand it over. \\
Refresh token & \fpath{httpOnly} cookie & 30 days. JavaScript cannot read it at all, so a cross-site scripting bug cannot steal it. \\
\bottomrule
\end{tabular}
\end{center}

You can see the cookie's exact settings by signing in with \fpath{curl}:

\begin{lstlisting}[language=bash]
curl -s -i -X POST http://localhost:8080/api/v1/auth/login \
  -H 'Content-Type: application/json' \
  -d '{"email":"responder@incidentflow.test","password":"incidentflow"}' \
  | grep -i set-cookie
\end{lstlisting}

\expect{A line containing \fpath{httponly}, \fpath{samesite=lax}, and
\fpath{path=/api/v1/auth}. That path is a deliberate narrowing: the cookie is
only ever sent to the four authentication endpoints, so it is not attached to
the hundreds of other API calls that have no use for it.}

\warnbox{In production the same cookie also carries \fpath{Secure}, which means
the browser will only send it over HTTPS. This has a consequence that has caught
many people, including during this project: deploy over plain HTTP and sign-in
\emph{appears} to work, then the session dies about fifteen minutes later when
the access token expires and the refresh cookie is never sent back. Nothing
errors, and no log records it. Chapter~\ref{ch:bugs-deploy} covers it.}

\section{Check your understanding}

\begin{itemize}[leftmargin=1.4cm]
\cbox Why does the sign-in error not tell you whether the email exists?
\cbox Why is the access token kept in memory rather than \fpath{localStorage}?
\cbox Why can JavaScript not read the refresh cookie?
\cbox What would break if \fpath{event.preventDefault()} were removed from
  \fpath{handleSubmit}?
\end{itemize}
